\documentclass[11pt]{article}
\usepackage[margin=1in]{geometry}
\usepackage{times}
\usepackage{hyperref}
\hypersetup{colorlinks=true, linkcolor=black, urlcolor=blue, citecolor=black}
\title{A Feasibility Study Is Not a Discovery: Boeing's GRASP Program and the Antigravity Rumor It Generated}
\author{Flyxion \\ \small Independent Researcher}
\date{}
\begin{document}
\maketitle

\begin{abstract}
Reports in 2002 that Boeing's Phantom Works division had a research effort, popularly referred to as GRASP, to investigate gravity modification circulated widely as evidence that a major aerospace contractor had found a credible path to gravity control propulsion. This paper argues that the available public record supports a considerably more modest reading. The acronym named a briefing document, Gravity Research for Advanced Space Propulsion, that proposed monitoring the claims of Podkletnov and conducting low-cost experiments, and a company statement quoted at the time said that no steps had been taken beyond that point. The persistence of the antigravity reading owes to the gap between what a large company's internal exploratory briefing actually is and what public reporting on leaked or rumored internal programs tends to imply. The paper is limited to what the public record shows; it cannot establish what Boeing did not do in private, and it says so.
\end{abstract}

\section{Introduction}

Trade press reporting in the summer of 2002, most notably by Nick Cook in \emph{Jane's Defence Weekly}, described a Boeing document referred to as GRASP, said to be investigating gravity modification technologies including some of the claims surveyed elsewhere in this series, above all Podkletnov's rotating superconductor work \cite{Cook2002}. The briefing obtained by the magazine stated that if gravity modification is real it would alter the entire aerospace business, and the reporting recorded that the head of the Phantom Works at the time acknowledged the company's interest in the work. The reporting, and much of the subsequent circulation of it, framed the existence of this document as evidence that a serious aerospace engineering organization considered gravity control technically credible, and by implication as third-party validation of the underlying antigravity claims themselves.

\section{What Large Engineering Organizations Actually Do With Speculative Claims}

Aerospace and defense contractors of Boeing's size maintain technology-scouting functions whose explicit purpose is to review emerging and even speculative claims across a wide range of fields, precisely so that a genuine breakthrough is not missed and so that competitors cannot claim exclusive awareness of a promising lead. The overwhelming majority of what such a function reviews does not pan out, and a review being conducted is evidence only that the claim was judged worth a bounded, low-cost look, not that the claim was judged likely to be true. Public reporting on internal corporate research efforts, especially when sourced from limited leaked documentation or from statements by researchers eager for their work to be taken seriously, systematically overstates the significance of a company having looked at something relative to a company having found something.

The same tendency affects the record of the principal claim. Public attention to the Podkletnov experiments in the 1990s came from a press story that the laboratory where he worked disavowed, and from a coauthor who stated that his name had been used without his prior knowledge, and Podkletnov himself withdrew a second paper after it had been accepted \cite{P2003}. A briefing that proposed to assess a claim with that history is a briefing that proposed to check it.

\section{What Actually Happened, So Far as the Public Record Shows}

A statement attributed to Boeing and circulated within weeks of the original story described the matter in terms that differ from the headlines. It said that GRASP is not a codename for a current project but an acronym for a presentation in which a Boeing engineer explains Podkletnov's theory and proposes that the company continue to monitor the work and perhaps conduct some low-cost experiments to assess its plausibility, and that no steps had been taken beyond that point \cite{Boeing2002}. The statement is a secondary quotation of a corporate response, and it should be read as such. It nonetheless fits the other features of the record: no published Boeing patent, technical paper, or flight demonstration has ever resulted from the effort, and no independent laboratory has reported a positive replication of the Podkletnov experiment to supply a basis for a larger program \cite{H2003,W2001}.

The original version of this essay described the effort as a program that was discontinued. The public record does not show that. It shows a briefing, a press story, a response, and silence. The silence is the expected outcome of a briefing that recommended monitoring a claim that was then tested elsewhere and not confirmed, and a company that does not pursue a claim does not announce that it is not pursuing it. This is a smaller conclusion than discontinuation, and it is the conclusion that the evidence supports.

\section{The Entropic Gravity Framing}

Nothing in the GRASP document's reported scope, insofar as the public record describes it, proposed a novel mechanism for gravitational coupling beyond the ones already addressed in this series, primarily the Podkletnov superconductor claim. A corporate feasibility review evaluating a mechanism that has no available channel to affect gravitational curvature, on either the standard general-relativistic or the entropic account, would be expected to conclude exactly what the available record suggests, regardless of how the review was initially reported.

\section{Objections}

It might be argued that a company would have strong incentive to keep a genuinely successful gravity modification result confidential for competitive and national security reasons, and that the absence of published results is therefore uninformative. This is a general-purpose objection available to defend any claim of suppressed success, and it proves too much: it would equally excuse the total absence of results for every device surveyed in this series, converting a complete lack of evidence into an argument for secrecy rather than for absence, a move that abandons ordinary evidentiary standards rather than satisfying them. The objection also cannot be answered from the public record, and this paper does not pretend to answer it. It observes only that an unfalsifiable account of what happened in private has no bearing on what the public record shows.

A second objection is that the briefing document, as reported, described the Podkletnov work in language more favorable than a skeptic would use, and that this indicates the engineer who wrote it believed the claim. It indicates that an internal advocate for investigating a technology wrote a document advocating investigation. Briefings that propose studies do not read as dismissals.

\section{Conclusion}

The public record surrounding Boeing's GRASP supports a mundane, if slightly deflating, conclusion: a briefing proposed to monitor and perhaps test a speculative claim, the claim was tested elsewhere without confirmation, and nothing followed that the public can see, a sequence of events considerably less remarkable than its circulation as an antigravity validation story has suggested. The paper cannot rule out undisclosed activity and does not rely on doing so.

\begin{thebibliography}{99}

\bibitem{Cook2002} N. Cook, report on Boeing's GRASP briefing and anti-gravity research, \emph{Jane's Defence Weekly}, 29 July 2002.

\bibitem{P2003} E. Podkletnov and G. Modanese, Impulse Gravity Generator Based on Charged YBa$_2$Cu$_3$O$_{7-y}$ Superconductor with Composite Crystal Structure, arXiv preprint, 2001, with the history of the withdrawn second paper summarized in the Wikipedia article on Eugene Podkletnov.

\bibitem{Boeing2002} Statement attributed to Boeing and quoted in discussion of the Jane's report, Extropians mailing list archive, August 2002.

\bibitem{H2003} G. Hathaway, B. Cleveland and Y. Bao, Gravity Modification Experiment Using a Rotating Superconducting Disk and Radio Frequency Fields, \emph{Physica C} 385 (2003) 488--500.

\bibitem{W2001} R. C. Woods, S. G. Cooke, J. Helme and C. H. Caldwell, Gravity Modification by High-Temperature Superconductors, AIAA Paper 2001-3363, 2001.

\end{thebibliography}

\end{document}
